The lower bound uses smooth controls. The upper argument must also handle
controls completed in the Gaussian covariance norm, where pointwise formulas
need not remain available. We need a precise relation between such controls, the
deterministic means they induce, and the wave operators generated by those
means. Spatial energy estimates provide the required bounds.

Recall the nonnegative-control variational value from
Section~\ref{S:unmasked-variational} and the wave phase from
Section~\ref{S:constant-phase-young}. We first compute a temporal equilibrium
density. Section~\ref{S:spatial-occupation-foundation} then extends physical
spatial pairings and defines the class of time-dependent spatial profiles used
in the later energy bounds. Section~\ref{S:actual-cameron-mean} constructs the
compact map from a Gaussian control to its mean on time blocks. Finally,
Section~\ref{S:young-upper-foundation} estimates perturbations of Young
evolution while retaining the deterministic endpoint. Together these results
prepare the Gaussian upper comparison. The spatial radius is still fixed; the
next chapter determines which bounds remain valid as it grows.
